\section{Introduction}
\label{sec:introduction}

Encrypted traffic classification has a reproducibility problem that is easy to
state. Papers report macro-F1 and accuracy above 95\%, sometimes above 99\%,
on datasets assembled from packet captures. Those numbers are real in the sense
that the code produces them. They are not results about traffic analysis.

We demonstrate this directly. Table~\ref{tab:split-contrast} reports the same
model, the same code, the same hyperparameters and the same data, evaluated
under four ways of partitioning the flows into training and test sets. On the
5G Traffic corpus the model scores macro-F1 \textbf{0.992} under a stratified
random split over flows --- the protocol used, to our reading, by the large
majority of published work in this area --- and \textbf{0.567} when flows from
one packet capture are constrained to a single side of the split. A 2016
statistical baseline behaves identically: 0.996 and 0.640.

The obvious diagnosis, that random splitting inflates results, is wrong, or at
least incomplete. On CESNET-QUIC22, a month of QUIC traffic from a national
backbone, we evaluated five split protocols at eight seeds each. Random,
temporal and client-disjoint splitting are statistically indistinguishable from
session-disjoint --- $p > 0.5$ and $|d| < 0.35$ in all three cases, which are
powered nulls rather than absences of evidence. What costs
\textbf{0.192 macro-F1} there is \emph{server}-disjoint splitting: testing only
on flows to servers never seen in training ($d = 17.0$, $p = 0.0078$, surviving
Holm-Bonferroni correction across the family).

\subsection*{The axis that matters is corpus-specific}

The two corpora fail in different places because they were collected
differently. 5G Traffic is a set of single-application packet captures: one
recording, one app, one device. A flow's capture is therefore nearly its label,
and a classifier that identifies the capture scores as well as one that
identifies the application. CESNET-QUIC22 is backbone traffic from many hosts
to many servers; its flows are already independent across days, so grouping by
day buys nothing, but the label correlates strongly with the destination, so
holding out servers bites.

Neither protocol protects against the other's failure. A paper reporting
session-disjoint results on backbone traffic has controlled nothing. A paper
reporting server-disjoint results on single-session captures has controlled
nothing either.

\subsection*{A diagnostic that runs before the model}

The right axis can be identified without training anything. We fit a decision
tree on a \emph{single} identifier column and report its macro-F1 under each
candidate protocol. On 5G Traffic the server address alone reaches 0.999 under
a random split and the capture-session identifier alone reaches 0.983; under
session-disjoint splitting the latter falls to exactly 0.000, which is the
grouping doing precisely what it claims. On CESNET-QUIC22 the SNI string alone
reaches 0.967 under every split --- unsurprising, since that dataset's labels
are derived from SNI, and it therefore measures the ceiling of the labelling
process rather than of any classifier.

These probes cost seconds. We argue they should precede any reported number in
this area, and we release them as a package that runs on any corpus in our
schema.

\subsection*{What survives under strict evaluation}

Having established the protocol, we use it. We evaluate a metric-trained
flow-level embedding against eight classical and four deep baselines on
identical committed splits, and we report five of our own architectural claims
failing their own pre-registered tests. Most notably, distance-based rejection
of unknown applications beats softmax thresholding by a wide margin --- 0.34
against 0.61 false-accept rate at matched true-positive rate --- but ablating
\emph{every} metric objective from the training makes rejection \emph{better},
not worse. The advantage belongs to the rejection rule, not to the
representation learning we introduced to produce it. We report this because a
reviewer would otherwise find it, and because the corrected claim is more
useful: it applies to anyone's embedding.

\subsection*{Contributions}

\begin{enumerate}
\item A measurement of how much the split protocol determines the reported
      result, on two corpora with opposite structure, every method on identical
      committed splits and every protocol at eight seeds
      (\S\ref{sec:protocols}). The same change is worth $0.426$ macro-F1 on one
      corpus and $0.001$ on the other. On the second, four of five axes are
      statistically null and the fifth costs $0.192$; the positive results
      survive Holm-Bonferroni correction and the nulls are powered rather than
      undetermined.
\item Single-identifier probes as a pre-training diagnostic for which axis a
      corpus requires, with the mechanism made explicit
      (\S\ref{sec:protocols}).
\item An evaluation of thirteen methods under strict protocols
      (\S\ref{sec:results}), including the finding that on a volume-defined
      taxonomy four independent neural architectures fall within $0.05$ macro-F1
      of each other while gradient-boosted trees on thirteen flow statistics sit
      $0.23$ above all of them, and an early-decision curve
      (Figure~\ref{fig:early}) and countermeasure-overhead evaluation
      (Figure~\ref{fig:robustness}) under the same protocols.
\item Five negative architectural results, each with the test that produced it
      (\S\ref{sec:negative}), including a representation-probing protocol
      showing an adversarial component leaving its nuisance as decodable as the
      raw input at every weight (Figure~\ref{fig:adversarial}), and a geometric
      explanation for why ablating our metric objectives \emph{improves}
      open-set rejection.
\item A sixth claim corrected rather than withdrawn: enrollment is flat within
      a corpus and recovers 91\% of a failed zero-shot transfer across corpora
      (\S\ref{sec:results}, Figure~\ref{fig:enrollment}).
\item An artifact containing every split manifest, every result including the
      failed experiments, the audit that found them, and a two-minute pipeline
      verification requiring no data download (\S\ref{sec:threats}).
\end{enumerate}

\subsection*{What we do not claim}

We do not claim state-of-the-art accuracy. On the corpus where our model is
strongest it is within one seed's noise of gradient-boosted trees on twenty
packet sizes, which train in three seconds. We did not run ET-BERT or the other
pre-trained byte-level transformers, because they tokenise payload bytes that
QUIC does not expose and our records do not retain; we say so rather than
citing their published numbers alongside ours.
